% ============================================================
%  TEACHING AGENT — Report Section
%  Included via % ============================================================
%  TEACHING AGENT — Report Section
%  Included via % ============================================================
%  TEACHING AGENT — Report Section
%  Included via % ============================================================
%  TEACHING AGENT — Report Section
%  Included via \input{teaching_agent} from project_report.tex
%
%  Required packages (already loaded in project_report.tex preamble):
%    tikz, booktabs, tabularx, xcolor, enumitem
% ============================================================

% ── Colour palette (safe to redefine if your report uses different names) ──
\definecolor{agentblue}{HTML}{2E75B6}
\definecolor{lighblue}{HTML}{D5E8F0}
\definecolor{agentgray}{HTML}{F2F2F2}
\definecolor{darkgray}{HTML}{404040}

% ============================================================
\section{Teaching Agent}
\label{app:teaching}
% ============================================================

% ──────────────────────────────────────────────────────────────
\subsection{Limitations and Future Scope}
\label{sec:ta-future}
% ──────────────────────────────────────────────────────────────

\subsubsection{Current Limitations}

\begin{enumerate}[leftmargin=*, label=\textbf{L\arabic*.}]

  \item \textbf{Structural-only Mermaid validation.}
  \texttt{validate\_mermaid()} checks for a recognised header and at least one content segment. A diagram that is structurally valid but semantically broken (e.g.\ malformed edge syntax that a renderer would reject) can pass this check. A full Mermaid parser is deferred to a future iteration.

  \item \textbf{Process-local conversation store.}
  \texttt{ConversationStore} lives in memory inside the worker process. A process restart clears all in-flight session histories. This is an accepted trade-off for a short-session tutoring product (sessions expire after 5 minutes of idle time), but would need to be replaced with a persistent store (Redis, database) for longer-lived sessions.

  \item \textbf{No cross-worker history sharing.}
  If the Teaching Agent is scaled horizontally, each worker process owns an independent \texttt{ConversationStore}. Requests from the same \texttt{sid} reaching different workers will not see each other's history. Sticky routing or an external shared store would be required for multi-process deployments.

  \item \textbf{Reflection quality is probabilistic.}
  The critique-revision loop is expected to improve output quality in the majority of runs, but it is a probabilistic guarantee. A cycle may occasionally produce a lateral move rather than an improvement. The automated test suite verifies control flow and token accounting; output quality validation requires real LLM calls and manual review.

\end{enumerate}

\subsubsection{Future Enhancements}

\begin{enumerate}[leftmargin=*, label=\textbf{F\arabic*.}]

  \item \textbf{Full Mermaid parsing.}
  Replace regex-based structural validation with a proper Mermaid parser to catch semantic errors before they reach the frontend renderer.

  \item \textbf{Persistent session storage.}
  Migrate \texttt{ConversationStore} to a Redis-backed or database-backed store to support multi-process deployments and survive worker restarts.

  \item \textbf{Adaptive mode selection.}
  Allow the agent (or the Planner) to adjust the \texttt{user\_level} dynamically based on learner performance signals from the Evaluation Agent, personalising the explanation depth over time.

  \item \textbf{Richer output formats.}
  Extend the four-section schema with optional sections such as \textit{Quiz Preview} or \textit{Related Topics} to provide tighter integration with downstream agents without breaking the existing contract.

  \item \textbf{Streaming reflection.}
  Stream revised content from the reflection loop back to the frontend in real time rather than waiting for all critique-revision cycles to complete before emitting the final output.

\end{enumerate}

 from project_report.tex
%
%  Required packages (already loaded in project_report.tex preamble):
%    tikz, booktabs, tabularx, xcolor, enumitem
% ============================================================

% ── Colour palette (safe to redefine if your report uses different names) ──
\definecolor{agentblue}{HTML}{2E75B6}
\definecolor{lighblue}{HTML}{D5E8F0}
\definecolor{agentgray}{HTML}{F2F2F2}
\definecolor{darkgray}{HTML}{404040}

% ============================================================
\section{Teaching Agent}
\label{app:teaching}
% ============================================================

% ──────────────────────────────────────────────────────────────
\subsection{Limitations and Future Scope}
\label{sec:ta-future}
% ──────────────────────────────────────────────────────────────

\subsubsection{Current Limitations}

\begin{enumerate}[leftmargin=*, label=\textbf{L\arabic*.}]

  \item \textbf{Structural-only Mermaid validation.}
  \texttt{validate\_mermaid()} checks for a recognised header and at least one content segment. A diagram that is structurally valid but semantically broken (e.g.\ malformed edge syntax that a renderer would reject) can pass this check. A full Mermaid parser is deferred to a future iteration.

  \item \textbf{Process-local conversation store.}
  \texttt{ConversationStore} lives in memory inside the worker process. A process restart clears all in-flight session histories. This is an accepted trade-off for a short-session tutoring product (sessions expire after 5 minutes of idle time), but would need to be replaced with a persistent store (Redis, database) for longer-lived sessions.

  \item \textbf{No cross-worker history sharing.}
  If the Teaching Agent is scaled horizontally, each worker process owns an independent \texttt{ConversationStore}. Requests from the same \texttt{sid} reaching different workers will not see each other's history. Sticky routing or an external shared store would be required for multi-process deployments.

  \item \textbf{Reflection quality is probabilistic.}
  The critique-revision loop is expected to improve output quality in the majority of runs, but it is a probabilistic guarantee. A cycle may occasionally produce a lateral move rather than an improvement. The automated test suite verifies control flow and token accounting; output quality validation requires real LLM calls and manual review.

\end{enumerate}

\subsubsection{Future Enhancements}

\begin{enumerate}[leftmargin=*, label=\textbf{F\arabic*.}]

  \item \textbf{Full Mermaid parsing.}
  Replace regex-based structural validation with a proper Mermaid parser to catch semantic errors before they reach the frontend renderer.

  \item \textbf{Persistent session storage.}
  Migrate \texttt{ConversationStore} to a Redis-backed or database-backed store to support multi-process deployments and survive worker restarts.

  \item \textbf{Adaptive mode selection.}
  Allow the agent (or the Planner) to adjust the \texttt{user\_level} dynamically based on learner performance signals from the Evaluation Agent, personalising the explanation depth over time.

  \item \textbf{Richer output formats.}
  Extend the four-section schema with optional sections such as \textit{Quiz Preview} or \textit{Related Topics} to provide tighter integration with downstream agents without breaking the existing contract.

  \item \textbf{Streaming reflection.}
  Stream revised content from the reflection loop back to the frontend in real time rather than waiting for all critique-revision cycles to complete before emitting the final output.

\end{enumerate}

 from project_report.tex
%
%  Required packages (already loaded in project_report.tex preamble):
%    tikz, booktabs, tabularx, xcolor, enumitem
% ============================================================

% ── Colour palette (safe to redefine if your report uses different names) ──
\definecolor{agentblue}{HTML}{2E75B6}
\definecolor{lighblue}{HTML}{D5E8F0}
\definecolor{agentgray}{HTML}{F2F2F2}
\definecolor{darkgray}{HTML}{404040}

% ============================================================
\section{Teaching Agent}
\label{app:teaching}
% ============================================================

% ──────────────────────────────────────────────────────────────
\subsection{Limitations and Future Scope}
\label{sec:ta-future}
% ──────────────────────────────────────────────────────────────

\subsubsection{Current Limitations}

\begin{enumerate}[leftmargin=*, label=\textbf{L\arabic*.}]

  \item \textbf{Structural-only Mermaid validation.}
  \texttt{validate\_mermaid()} checks for a recognised header and at least one content segment. A diagram that is structurally valid but semantically broken (e.g.\ malformed edge syntax that a renderer would reject) can pass this check. A full Mermaid parser is deferred to a future iteration.

  \item \textbf{Process-local conversation store.}
  \texttt{ConversationStore} lives in memory inside the worker process. A process restart clears all in-flight session histories. This is an accepted trade-off for a short-session tutoring product (sessions expire after 5 minutes of idle time), but would need to be replaced with a persistent store (Redis, database) for longer-lived sessions.

  \item \textbf{No cross-worker history sharing.}
  If the Teaching Agent is scaled horizontally, each worker process owns an independent \texttt{ConversationStore}. Requests from the same \texttt{sid} reaching different workers will not see each other's history. Sticky routing or an external shared store would be required for multi-process deployments.

  \item \textbf{Reflection quality is probabilistic.}
  The critique-revision loop is expected to improve output quality in the majority of runs, but it is a probabilistic guarantee. A cycle may occasionally produce a lateral move rather than an improvement. The automated test suite verifies control flow and token accounting; output quality validation requires real LLM calls and manual review.

\end{enumerate}

\subsubsection{Future Enhancements}

\begin{enumerate}[leftmargin=*, label=\textbf{F\arabic*.}]

  \item \textbf{Full Mermaid parsing.}
  Replace regex-based structural validation with a proper Mermaid parser to catch semantic errors before they reach the frontend renderer.

  \item \textbf{Persistent session storage.}
  Migrate \texttt{ConversationStore} to a Redis-backed or database-backed store to support multi-process deployments and survive worker restarts.

  \item \textbf{Adaptive mode selection.}
  Allow the agent (or the Planner) to adjust the \texttt{user\_level} dynamically based on learner performance signals from the Evaluation Agent, personalising the explanation depth over time.

  \item \textbf{Richer output formats.}
  Extend the four-section schema with optional sections such as \textit{Quiz Preview} or \textit{Related Topics} to provide tighter integration with downstream agents without breaking the existing contract.

  \item \textbf{Streaming reflection.}
  Stream revised content from the reflection loop back to the frontend in real time rather than waiting for all critique-revision cycles to complete before emitting the final output.

\end{enumerate}

 from project_report.tex
%
%  Required packages (already loaded in project_report.tex preamble):
%    tikz, booktabs, tabularx, xcolor, enumitem
% ============================================================

% ── Colour palette (safe to redefine if your report uses different names) ──
\definecolor{agentblue}{HTML}{2E75B6}
\definecolor{lighblue}{HTML}{D5E8F0}
\definecolor{agentgray}{HTML}{F2F2F2}
\definecolor{darkgray}{HTML}{404040}

% ============================================================
\section{Teaching Agent}
\label{app:teaching}
% ============================================================

% ──────────────────────────────────────────────────────────────
\subsection{Limitations and Future Scope}
\label{sec:ta-future}
% ──────────────────────────────────────────────────────────────

\subsubsection{Current Limitations}

\begin{enumerate}[leftmargin=*, label=\textbf{L\arabic*.}]

  \item \textbf{Structural-only Mermaid validation.}
  \texttt{validate\_mermaid()} checks for a recognised header and at least one content segment. A diagram that is structurally valid but semantically broken (e.g.\ malformed edge syntax that a renderer would reject) can pass this check. A full Mermaid parser is deferred to a future iteration.

  \item \textbf{Process-local conversation store.}
  \texttt{ConversationStore} lives in memory inside the worker process. A process restart clears all in-flight session histories. This is an accepted trade-off for a short-session tutoring product (sessions expire after 5 minutes of idle time), but would need to be replaced with a persistent store (Redis, database) for longer-lived sessions.

  \item \textbf{No cross-worker history sharing.}
  If the Teaching Agent is scaled horizontally, each worker process owns an independent \texttt{ConversationStore}. Requests from the same \texttt{sid} reaching different workers will not see each other's history. Sticky routing or an external shared store would be required for multi-process deployments.

  \item \textbf{Reflection quality is probabilistic.}
  The critique-revision loop is expected to improve output quality in the majority of runs, but it is a probabilistic guarantee. A cycle may occasionally produce a lateral move rather than an improvement. The automated test suite verifies control flow and token accounting; output quality validation requires real LLM calls and manual review.

\end{enumerate}

\subsubsection{Future Enhancements}

\begin{enumerate}[leftmargin=*, label=\textbf{F\arabic*.}]

  \item \textbf{Full Mermaid parsing.}
  Replace regex-based structural validation with a proper Mermaid parser to catch semantic errors before they reach the frontend renderer.

  \item \textbf{Persistent session storage.}
  Migrate \texttt{ConversationStore} to a Redis-backed or database-backed store to support multi-process deployments and survive worker restarts.

  \item \textbf{Adaptive mode selection.}
  Allow the agent (or the Planner) to adjust the \texttt{user\_level} dynamically based on learner performance signals from the Evaluation Agent, personalising the explanation depth over time.

  \item \textbf{Richer output formats.}
  Extend the four-section schema with optional sections such as \textit{Quiz Preview} or \textit{Related Topics} to provide tighter integration with downstream agents without breaking the existing contract.

  \item \textbf{Streaming reflection.}
  Stream revised content from the reflection loop back to the frontend in real time rather than waiting for all critique-revision cycles to complete before emitting the final output.

\end{enumerate}

